% --- Archon physics formalization source begin ---
% archon:physics
% archon:covers PhyXMiniProblems/problem_phyx_mini_0124.lean
% archon:source-report reports/phyx_mini/problem_phyx_mini_0124.source.json

\chapter{Physics problem phyx\_mini\_0124}
\label{ch:PhyXMiniProblems_problem_phyx_mini_0124}

\paragraph{Problem source.}
\#\# Physical scenario

An opaque barrier has an inner membrane and an outer membrane that slide past each other, as shown in figure. Each membrane includes parallel slits of width \$a\$ separated by a distance \$d\$. A screen forms a circular arc subtending \$60\textasciicircum{}\textbackslash{}circ\$ at the fixed midpoint between the slits. A green 532 nm laser impinges on the slits from the left. The outer membrane moves upward with speed \$v\$ while the inner membrane moves downward with the same speed, propelled by nanomotors. At time \$t = 0\$, point \$P\$ on the outer membrane is adjacent to point \$Q\$ on the inner membrane so that the effective aperture width is zero. The aperture is fully closed again at \$t = 3.00\$ s.

\#\# Question

What are the angular positions of these spots?

\#\# Answer choices

- A: A: 2.88°

- B: B: 2.48°

- C: C: 2.58°

- D: D: 2.68°

\#\# Figure caption (auxiliary; use the image as primary evidence)

The image is a diagram illustrating a two-slit diffraction experiment setup. Here's a detailed description:

1. **Objects and Layout**:
   - There is a barrier with two slits, marked by the vertical gray lines. The slits are labeled \textbackslash{}( P \textbackslash{}) and \textbackslash{}( Q \textbackslash{}).
   - A screen is on the right side, which is curved, suggesting a focus on detecting interference patterns.
   - Blue arcs emanate from the slits, representing wavefronts or wave patterns from each slit.

2. **Text and Labels**:
   - Distance labels \textbackslash{}( d \textbackslash{}) and \textbackslash{}( a \textbackslash{}) indicate measurements between points. \textbackslash{}( d \textbackslash{}) is the distance between the slits, and \textbackslash{}( a \textbackslash{}) is from the slit Q downwards to another point.
   - There is an angle marked \textbackslash{}( \textbackslash{}theta \textbackslash{}) (theta) from the line between the slits and a line to the screen.
   - The text "Not to scale" is written at the bottom of the diagram.
   - Velocity vectors \textbackslash{}( \textbackslash{}vec\{v\} \textbackslash{}) are shown as green arrows moving through the slits.
  
3. **Relationships and Indications**:
   - The design suggests wave interference patterns, commonly associated with the double-slit experiment.
   - The angle \textbackslash{}( \textbackslash{}theta \textbackslash{}) likely indicates the angle from the central axis of the pattern to a particular point on the screen.
   - The green arrows show the direction of particle or wave movement towards the slits.

Overall, the diagram illustrates key aspects of wave interference seen in a double-slit experiment.

\#\# Dataset metadata

Wave Optics; Physical Model Grounding Reasoning, Multi-Formula Reasoning

\paragraph{Recorded answer/context.}
A: A: 2.88°

\paragraph{Figure/image path.}
/root/proposal\_for\_physic/science-mango/phyx\_mini\_run/phyx\_data/test\_image/124.png

\paragraph{Formalization target.}
create a compiling Lean file with sorry bodies at `PhyXMiniProblems/problem\_phyx\_mini\_0124.lean`. The Lean declarations must preserve the physical quantities, dimensions or dimensional roles, figure labels, governing-law hypotheses, and final relation expressed by this problem.
Use Mathlib/Physlib names found through LeanExplore where available. If a physics API is missing, introduce faithful local abstractions rather than scalar placeholder aliases.

\begin{theorem}[Physics formalization target]
\label{thm:physics:phyx_mini_0124:target}
\lean{PhyXMiniProblems.ProblemPhyXMini0124.problem_phyx_mini_0124}
\uses{def:physics:phyx-mini-0124:phyxminiproblems-problemphyxmini0124-lengthinmeters, def:physics:phyx-mini-0124:phyxminiproblems-problemphyxmini0124-speedinmeterspersecond, def:physics:phyx-mini-0124:phyxminiproblems-problemphyxmini0124-slidingmembraneinterferencesetup, def:physics:phyx-mini-0124:phyxminiproblems-problemphyxmini0124-spotangle, def:physics:phyx-mini-0124:phyxminiproblems-problemphyxmini0124-matchesfiguregeometry, def:physics:phyx-mini-0124:phyxminiproblems-problemphyxmini0124-matchesproblemreadouts, def:physics:phyx-mini-0124:phyxminiproblems-problemphyxmini0124-hasphysicalparameters, def:physics:phyx-mini-0124:phyxminiproblems-problemphyxmini0124-satisfiesslidingmembranekinematics, def:physics:phyx-mini-0124:phyxminiproblems-problemphyxmini0124-satisfiesslitarrayinterferencelaw}
The assigned autoformalize agent should translate this physics problem into Lean declarations in the covered file, with theorem and lemma proof bodies written as `by sorry`.
\end{theorem}
\begin{proof}
This is an autoformalization task, not a proof task. Produce statements that can later be proved by the physics prover without weakening the physical model.
\end{proof}

% --- Archon declaration topology begin ---
\section{Declaration topology}
\begin{definition}[Length Quantity]
  \label{def:physics:phyx-mini-0124:phyxminiproblems-problemphyxmini0124-lengthquantity}
  \lean{PhyXMiniProblems.ProblemPhyXMini0124.LengthQuantity}
  A nonnegative unit-independent physical length
\end{definition}
\begin{proof}
  This is the declared model definition of Length Quantity.
\end{proof}
\begin{definition}[Signed Length Quantity]
  \label{def:physics:phyx-mini-0124:phyxminiproblems-problemphyxmini0124-signedlengthquantity}
  \lean{PhyXMiniProblems.ProblemPhyXMini0124.SignedLengthQuantity}
  A signed transverse coordinate with the physical dimension of length
\end{definition}
\begin{proof}
  This is the declared model definition of Signed Length Quantity.
\end{proof}
\begin{definition}[Time Quantity]
  \label{def:physics:phyx-mini-0124:phyxminiproblems-problemphyxmini0124-timequantity}
  \lean{PhyXMiniProblems.ProblemPhyXMini0124.TimeQuantity}
  A nonnegative unit-independent physical time
\end{definition}
\begin{proof}
  This is the declared model definition of Time Quantity.
\end{proof}
\begin{definition}[Speed Quantity]
  \label{def:physics:phyx-mini-0124:phyxminiproblems-problemphyxmini0124-speedquantity}
  \lean{PhyXMiniProblems.ProblemPhyXMini0124.SpeedQuantity}
  The nonnegative physical speed type supplied by Physlib
\end{definition}
\begin{proof}
  This is the declared model definition of Speed Quantity.
\end{proof}
\begin{definition}[length Readout]
  \label{def:physics:phyx-mini-0124:phyxminiproblems-problemphyxmini0124-lengthreadout}
  \lean{PhyXMiniProblems.ProblemPhyXMini0124.lengthReadout}
  \uses{def:physics:phyx-mini-0124:phyxminiproblems-problemphyxmini0124-lengthquantity}
  Read a nonnegative physical length in a selected length unit
\end{definition}
\begin{proof}
  This is the declared model definition of length Readout.
\end{proof}
\begin{definition}[signed Length Readout]
  \label{def:physics:phyx-mini-0124:phyxminiproblems-problemphyxmini0124-signedlengthreadout}
  \lean{PhyXMiniProblems.ProblemPhyXMini0124.signedLengthReadout}
  \uses{def:physics:phyx-mini-0124:phyxminiproblems-problemphyxmini0124-signedlengthquantity}
  Read a signed physical position in a selected length unit
\end{definition}
\begin{proof}
  This is the declared model definition of signed Length Readout.
\end{proof}
\begin{definition}[time Readout]
  \label{def:physics:phyx-mini-0124:phyxminiproblems-problemphyxmini0124-timereadout}
  \lean{PhyXMiniProblems.ProblemPhyXMini0124.timeReadout}
  \uses{def:physics:phyx-mini-0124:phyxminiproblems-problemphyxmini0124-timequantity}
  Read a physical time in a selected time unit
\end{definition}
\begin{proof}
  This is the declared model definition of time Readout.
\end{proof}
\begin{definition}[speed Readout]
  \label{def:physics:phyx-mini-0124:phyxminiproblems-problemphyxmini0124-speedreadout}
  \lean{PhyXMiniProblems.ProblemPhyXMini0124.speedReadout}
  \uses{def:physics:phyx-mini-0124:phyxminiproblems-problemphyxmini0124-speedquantity}
  Read a physical speed in selected length and time units
\end{definition}
\begin{proof}
  This is the declared model definition of speed Readout.
\end{proof}
\begin{definition}[length In Nanometers]
  \label{def:physics:phyx-mini-0124:phyxminiproblems-problemphyxmini0124-lengthinnanometers}
  \lean{PhyXMiniProblems.ProblemPhyXMini0124.lengthInNanometers}
  \uses{def:physics:phyx-mini-0124:phyxminiproblems-problemphyxmini0124-lengthquantity, def:physics:phyx-mini-0124:phyxminiproblems-problemphyxmini0124-lengthreadout}
  The nanometer readout used for the laser wavelength
\end{definition}
\begin{proof}
  This is the declared model definition of length In Nanometers.
\end{proof}
\begin{definition}[length In Meters]
  \label{def:physics:phyx-mini-0124:phyxminiproblems-problemphyxmini0124-lengthinmeters}
  \lean{PhyXMiniProblems.ProblemPhyXMini0124.lengthInMeters}
  \uses{def:physics:phyx-mini-0124:phyxminiproblems-problemphyxmini0124-lengthquantity, def:physics:phyx-mini-0124:phyxminiproblems-problemphyxmini0124-lengthreadout}
  The meter readout used by the kinematic and interference laws
\end{definition}
\begin{proof}
  This is the declared model definition of length In Meters.
\end{proof}
\begin{definition}[time In Seconds]
  \label{def:physics:phyx-mini-0124:phyxminiproblems-problemphyxmini0124-timeinseconds}
  \lean{PhyXMiniProblems.ProblemPhyXMini0124.timeInSeconds}
  \uses{def:physics:phyx-mini-0124:phyxminiproblems-problemphyxmini0124-timequantity, def:physics:phyx-mini-0124:phyxminiproblems-problemphyxmini0124-timereadout}
  The second readout used for the closure cycle
\end{definition}
\begin{proof}
  This is the declared model definition of time In Seconds.
\end{proof}
\begin{definition}[speed In Meters Per Second]
  \label{def:physics:phyx-mini-0124:phyxminiproblems-problemphyxmini0124-speedinmeterspersecond}
  \lean{PhyXMiniProblems.ProblemPhyXMini0124.speedInMetersPerSecond}
  \uses{def:physics:phyx-mini-0124:phyxminiproblems-problemphyxmini0124-speedquantity, def:physics:phyx-mini-0124:phyxminiproblems-problemphyxmini0124-speedreadout}
  The SI speed readout in meters per second
\end{definition}
\begin{proof}
  This is the declared model definition of speed In Meters Per Second.
\end{proof}
\begin{definition}[Membrane Label]
  \label{def:physics:phyx-mini-0124:phyxminiproblems-problemphyxmini0124-membranelabel}
  \lean{PhyXMiniProblems.ProblemPhyXMini0124.MembraneLabel}
  The two independently moving layers of the opaque barrier
\end{definition}
\begin{proof}
  This is the declared model definition of Membrane Label.
\end{proof}
\begin{definition}[Vertical Direction]
  \label{def:physics:phyx-mini-0124:phyxminiproblems-problemphyxmini0124-verticaldirection}
  \lean{PhyXMiniProblems.ProblemPhyXMini0124.VerticalDirection}
  The two opposite transverse directions shown by the green arrows
\end{definition}
\begin{proof}
  This is the declared model definition of Vertical Direction.
\end{proof}
\begin{definition}[Figure Point Label]
  \label{def:physics:phyx-mini-0124:phyxminiproblems-problemphyxmini0124-figurepointlabel}
  \lean{PhyXMiniProblems.ProblemPhyXMini0124.FigurePointLabel}
  The two moving edge points labelled in the primary figure
\end{definition}
\begin{proof}
  This is the declared model definition of Figure Point Label.
\end{proof}
\begin{definition}[Barrier Kind]
  \label{def:physics:phyx-mini-0124:phyxminiproblems-problemphyxmini0124-barrierkind}
  \lean{PhyXMiniProblems.ProblemPhyXMini0124.BarrierKind}
  Qualitative material character stated for the barrier
\end{definition}
\begin{proof}
  This is the declared model definition of Barrier Kind.
\end{proof}
\begin{definition}[Drive Kind]
  \label{def:physics:phyx-mini-0124:phyxminiproblems-problemphyxmini0124-drivekind}
  \lean{PhyXMiniProblems.ProblemPhyXMini0124.DriveKind}
  The actuator named in the problem statement
\end{definition}
\begin{proof}
  This is the declared model definition of Drive Kind.
\end{proof}
\begin{definition}[Incident Side]
  \label{def:physics:phyx-mini-0124:phyxminiproblems-problemphyxmini0124-incidentside}
  \lean{PhyXMiniProblems.ProblemPhyXMini0124.IncidentSide}
  The side from which the beam reaches the membrane stack
\end{definition}
\begin{proof}
  This is the declared model definition of Incident Side.
\end{proof}
\begin{definition}[Light Color]
  \label{def:physics:phyx-mini-0124:phyxminiproblems-problemphyxmini0124-lightcolor}
  \lean{PhyXMiniProblems.ProblemPhyXMini0124.LightColor}
  The visible color stated for the 532 nm laser
\end{definition}
\begin{proof}
  This is the declared model definition of Light Color.
\end{proof}
\begin{definition}[Beam Kind]
  \label{def:physics:phyx-mini-0124:phyxminiproblems-problemphyxmini0124-beamkind}
  \lean{PhyXMiniProblems.ProblemPhyXMini0124.BeamKind}
  The coherent monochromatic beam character implicit in an ideal laser
\end{definition}
\begin{proof}
  This is the declared model definition of Beam Kind.
\end{proof}
\begin{definition}[Screen Shape]
  \label{def:physics:phyx-mini-0124:phyxminiproblems-problemphyxmini0124-screenshape}
  \lean{PhyXMiniProblems.ProblemPhyXMini0124.ScreenShape}
  The screen shape visible in the figure
\end{definition}
\begin{proof}
  This is the declared model definition of Screen Shape.
\end{proof}
\begin{definition}[Screen Center]
  \label{def:physics:phyx-mini-0124:phyxminiproblems-problemphyxmini0124-screencenter}
  \lean{PhyXMiniProblems.ProblemPhyXMini0124.ScreenCenter}
  The fixed point about which screen angles are measured
\end{definition}
\begin{proof}
  This is the declared model definition of Screen Center.
\end{proof}
\begin{definition}[Optical Regime]
  \label{def:physics:phyx-mini-0124:phyxminiproblems-problemphyxmini0124-opticalregime}
  \lean{PhyXMiniProblems.ProblemPhyXMini0124.OpticalRegime}
  The ideal optical regime in which the stated interference law applies
\end{definition}
\begin{proof}
  This is the declared model definition of Optical Regime.
\end{proof}
\begin{definition}[Spot Label]
  \label{def:physics:phyx-mini-0124:phyxminiproblems-problemphyxmini0124-spotlabel}
  \lean{PhyXMiniProblems.ProblemPhyXMini0124.SpotLabel}
  The two symmetric spots whose angular positions are requested
\end{definition}
\begin{proof}
  This is the declared model definition of Spot Label.
\end{proof}
\begin{definition}[spot Order]
  \label{def:physics:phyx-mini-0124:phyxminiproblems-problemphyxmini0124-spotorder}
  \lean{PhyXMiniProblems.ProblemPhyXMini0124.spotOrder}
  \uses{def:physics:phyx-mini-0124:phyxminiproblems-problemphyxmini0124-spotlabel}
  The signed diffraction/interference order assigned to each named spot
\end{definition}
\begin{proof}
  This is the declared model definition of spot Order.
\end{proof}
\begin{definition}[Sliding Membrane Interference Setup]
  \label{def:physics:phyx-mini-0124:phyxminiproblems-problemphyxmini0124-slidingmembraneinterferencesetup}
  \lean{PhyXMiniProblems.ProblemPhyXMini0124.SlidingMembraneInterferenceSetup}
  \uses{def:physics:phyx-mini-0124:phyxminiproblems-problemphyxmini0124-lengthquantity, def:physics:phyx-mini-0124:phyxminiproblems-problemphyxmini0124-signedlengthquantity, def:physics:phyx-mini-0124:phyxminiproblems-problemphyxmini0124-timequantity, def:physics:phyx-mini-0124:phyxminiproblems-problemphyxmini0124-speedquantity, def:physics:phyx-mini-0124:phyxminiproblems-problemphyxmini0124-membranelabel, def:physics:phyx-mini-0124:phyxminiproblems-problemphyxmini0124-verticaldirection, def:physics:phyx-mini-0124:phyxminiproblems-problemphyxmini0124-figurepointlabel, def:physics:phyx-mini-0124:phyxminiproblems-problemphyxmini0124-barrierkind, def:physics:phyx-mini-0124:phyxminiproblems-problemphyxmini0124-drivekind, def:physics:phyx-mini-0124:phyxminiproblems-problemphyxmini0124-incidentside, def:physics:phyx-mini-0124:phyxminiproblems-problemphyxmini0124-lightcolor, def:physics:phyx-mini-0124:phyxminiproblems-problemphyxmini0124-beamkind, def:physics:phyx-mini-0124:phyxminiproblems-problemphyxmini0124-screenshape, def:physics:phyx-mini-0124:phyxminiproblems-problemphyxmini0124-screencenter, def:physics:phyx-mini-0124:phyxminiproblems-problemphyxmini0124-opticalregime}
  All physical quantities and figure-indexed observations in the setup
\end{definition}
\begin{proof}
  This is the declared model definition of Sliding Membrane Interference Setup.
\end{proof}
\begin{definition}[spot Angle]
  \label{def:physics:phyx-mini-0124:phyxminiproblems-problemphyxmini0124-spotangle}
  \lean{PhyXMiniProblems.ProblemPhyXMini0124.spotAngle}
  \uses{def:physics:phyx-mini-0124:phyxminiproblems-problemphyxmini0124-spotlabel, def:physics:phyx-mini-0124:phyxminiproblems-problemphyxmini0124-spotorder, def:physics:phyx-mini-0124:phyxminiproblems-problemphyxmini0124-slidingmembraneinterferencesetup}
  The angular position of either named first-order spot
\end{definition}
\begin{proof}
  This is the declared model definition of spot Angle.
\end{proof}
\begin{definition}[Fully Closed At]
  \label{def:physics:phyx-mini-0124:phyxminiproblems-problemphyxmini0124-fullyclosedat}
  \lean{PhyXMiniProblems.ProblemPhyXMini0124.FullyClosedAt}
  \uses{def:physics:phyx-mini-0124:phyxminiproblems-problemphyxmini0124-timequantity, def:physics:phyx-mini-0124:phyxminiproblems-problemphyxmini0124-lengthinmeters, def:physics:phyx-mini-0124:phyxminiproblems-problemphyxmini0124-slidingmembraneinterferencesetup}
  The effective overlap aperture is fully closed at a given instant
\end{definition}
\begin{proof}
  This is the declared model definition of Fully Closed At.
\end{proof}
\begin{definition}[Within First Aperture Cycle]
  \label{def:physics:phyx-mini-0124:phyxminiproblems-problemphyxmini0124-withinfirstaperturecycle}
  \lean{PhyXMiniProblems.ProblemPhyXMini0124.WithinFirstApertureCycle}
  \uses{def:physics:phyx-mini-0124:phyxminiproblems-problemphyxmini0124-timequantity, def:physics:phyx-mini-0124:phyxminiproblems-problemphyxmini0124-timeinseconds, def:physics:phyx-mini-0124:phyxminiproblems-problemphyxmini0124-slidingmembraneinterferencesetup}
  A time belongs to the first closure-to-closure cycle
\end{definition}
\begin{proof}
  This is the declared model definition of Within First Aperture Cycle.
\end{proof}
\begin{definition}[Matches Figure Geometry]
  \label{def:physics:phyx-mini-0124:phyxminiproblems-problemphyxmini0124-matchesfiguregeometry}
  \lean{PhyXMiniProblems.ProblemPhyXMini0124.MatchesFigureGeometry}
  \uses{def:physics:phyx-mini-0124:phyxminiproblems-problemphyxmini0124-spotorder, def:physics:phyx-mini-0124:phyxminiproblems-problemphyxmini0124-slidingmembraneinterferencesetup, def:physics:phyx-mini-0124:phyxminiproblems-problemphyxmini0124-spotangle}
  Qualitative geometry read from the primary image. In particular, the angle signs and screen bounds select the principal representatives of the two first-order spots without assigning either spot a numerical answer value
\end{definition}
\begin{proof}
  This is the declared model definition of Matches Figure Geometry.
\end{proof}
\begin{definition}[Matches Problem Readouts]
  \label{def:physics:phyx-mini-0124:phyxminiproblems-problemphyxmini0124-matchesproblemreadouts}
  \lean{PhyXMiniProblems.ProblemPhyXMini0124.MatchesProblemReadouts}
  \uses{def:physics:phyx-mini-0124:phyxminiproblems-problemphyxmini0124-lengthinnanometers, def:physics:phyx-mini-0124:phyxminiproblems-problemphyxmini0124-timeinseconds, def:physics:phyx-mini-0124:phyxminiproblems-problemphyxmini0124-slidingmembraneinterferencesetup, def:physics:phyx-mini-0124:phyxminiproblems-problemphyxmini0124-fullyclosedat}
  Numerical and endpoint readouts explicitly supplied by the problem. This contains 532 nm, 0 s, 3 s, and the 60 degree screen span, but no spot angle or answer choice
\end{definition}
\begin{proof}
  This is the declared model definition of Matches Problem Readouts.
\end{proof}
\begin{definition}[Has Physical Parameters]
  \label{def:physics:phyx-mini-0124:phyxminiproblems-problemphyxmini0124-hasphysicalparameters}
  \lean{PhyXMiniProblems.ProblemPhyXMini0124.HasPhysicalParameters}
  \uses{def:physics:phyx-mini-0124:phyxminiproblems-problemphyxmini0124-lengthinmeters, def:physics:phyx-mini-0124:phyxminiproblems-problemphyxmini0124-timeinseconds, def:physics:phyx-mini-0124:phyxminiproblems-problemphyxmini0124-speedinmeterspersecond, def:physics:phyx-mini-0124:phyxminiproblems-problemphyxmini0124-slidingmembraneinterferencesetup, def:physics:phyx-mini-0124:phyxminiproblems-problemphyxmini0124-withinfirstaperturecycle}
  Positivity and nondegeneracy of the experiment. The existence of an open intermediate aperture rules out a barrier which remains closed throughout
\end{definition}
\begin{proof}
  This is the declared model definition of Has Physical Parameters.
\end{proof}
\begin{definition}[Satisfies Sliding Membrane Kinematics]
  \label{def:physics:phyx-mini-0124:phyxminiproblems-problemphyxmini0124-satisfiesslidingmembranekinematics}
  \lean{PhyXMiniProblems.ProblemPhyXMini0124.SatisfiesSlidingMembraneKinematics}
  \uses{def:physics:phyx-mini-0124:phyxminiproblems-problemphyxmini0124-lengthreadout, def:physics:phyx-mini-0124:phyxminiproblems-problemphyxmini0124-signedlengthreadout, def:physics:phyx-mini-0124:phyxminiproblems-problemphyxmini0124-speedreadout, def:physics:phyx-mini-0124:phyxminiproblems-problemphyxmini0124-timeinseconds, def:physics:phyx-mini-0124:phyxminiproblems-problemphyxmini0124-slidingmembraneinterferencesetup, def:physics:phyx-mini-0124:phyxminiproblems-problemphyxmini0124-withinfirstaperturecycle}
  Uniform opposite membrane motion and periodic reclosure. Between successive closures the relative displacement is one repeat pitch d, hence d = 2 v Δt. This general kinematic law contains no optical answer
\end{definition}
\begin{proof}
  This is the declared model definition of Satisfies Sliding Membrane Kinematics.
\end{proof}
\begin{definition}[Satisfies Slit Array Interference Law]
  \label{def:physics:phyx-mini-0124:phyxminiproblems-problemphyxmini0124-satisfiesslitarrayinterferencelaw}
  \lean{PhyXMiniProblems.ProblemPhyXMini0124.SatisfiesSlitArrayInterferenceLaw}
  \uses{def:physics:phyx-mini-0124:phyxminiproblems-problemphyxmini0124-lengthinmeters, def:physics:phyx-mini-0124:phyxminiproblems-problemphyxmini0124-slidingmembraneinterferencesetup}
  Far-field slit-array interference for every visible signed order: d sin(θₘ) = m λ. Visibility is explicit because orders outside the screen or diffraction envelope need not produce observed spots
\end{definition}
\begin{proof}
  This is the declared model definition of Satisfies Slit Array Interference Law.
\end{proof}
\begin{definition}[Answer Choice]
  \label{def:physics:phyx-mini-0124:phyxminiproblems-problemphyxmini0124-answerchoice}
  \lean{PhyXMiniProblems.ProblemPhyXMini0124.AnswerChoice}
  The four answer labels printed with the question
\end{definition}
\begin{proof}
  This is the declared model definition of Answer Choice.
\end{proof}
\begin{definition}[displayed Magnitude Degrees]
  \label{def:physics:phyx-mini-0124:phyxminiproblems-problemphyxmini0124-answerchoice-displayedmagnitudedegrees}
  \lean{PhyXMiniProblems.ProblemPhyXMini0124.AnswerChoice.displayedMagnitudeDegrees}
  \uses{def:physics:phyx-mini-0124:phyxminiproblems-problemphyxmini0124-answerchoice}
  The positive angular magnitude in degrees displayed beside each choice
\end{definition}
\begin{proof}
  This is the declared model definition of displayed Magnitude Degrees.
\end{proof}
\begin{definition}[recorded Answer Choice]
  \label{def:physics:phyx-mini-0124:phyxminiproblems-problemphyxmini0124-recordedanswerchoice}
  \lean{PhyXMiniProblems.ProblemPhyXMini0124.recordedAnswerChoice}
  \uses{def:physics:phyx-mini-0124:phyxminiproblems-problemphyxmini0124-answerchoice}
  The dataset records choice A; this is metadata, not a physics premise
\end{definition}
\begin{proof}
  This is the declared model definition of recorded Answer Choice.
\end{proof}
% --- Archon declaration topology end ---
% --- Archon physics formalization source end ---
